\section{Methodology}

The proposed \textsf{StructSynth} leverages a Large Language Model (LLM) to perform efficient dependency graph discovery and subsequent data generation. As depicted in Figure~\ref{fig:pipeline}, the methodology comprises two primary stages: a \textbf{structure learning phase} to discover the underlying dependency graph and a \textbf{graph-aware generation phase} that synthesizes new tabular data using the learned graph as a blueprint.\footnote{The complete algorithmic procedure is provided in Algorithm~\ref{alg:structsynth}. Detailed prompt templates with formatting examples are available in Appendix~\ref{app:prompts}.}

\subsection{Problem Definition}

Consider a small tabular training dataset, denoted as $\mathcal{D}_{\mathtt{train}} = (\mathbf{X}_{\mathtt{train}}, \mathbf{A})$. $\mathbf{X}_{\mathtt{train}} = \{\mathbf{x}_1, \ldots, \mathbf{x}_n\}$ represents a set of $n$ instances (or rows), where the number of available samples is severely limited ($n \le 100$). The set $\mathbf{A} = \{A_1, \ldots, A_K\}$ defines the schema of the data through its $K$ attributes (or columns). Our primary objective is to employ an LLM-based generation function, $f_{\mathtt{LLM}}$, to produce a synthetic dataset, $\mathcal{D}_{\mathtt{synth}} = (\mathbf{X}_{\mathtt{synth}}, \mathbf{A})$, that is both realistic and diverse. The ultimate goal is to demonstrate that a downstream model trained on an augmented dataset, $\mathcal{D}_{\mathtt{aug}} = (\mathbf{X}_{\mathtt{train}} \cup \mathbf{X}_{\mathtt{synth}}, \mathbf{A})$, exhibits superior performance on a held-out test set, $\mathcal{D}_{\mathtt{test}}$, when compared to a model trained exclusively on the original, smaller dataset $\mathcal{D}_{\mathtt{train}}$.

\subsection{Structural Formulation: Why a DAG?}
\label{sec:why_dag}

A central design choice in \textsf{StructSynth} is the representation of inter-feature dependencies as a Directed Acyclic Graph (DAG). Since tabular feature relationships are often associative rather than inherently directional, this choice warrants explicit justification. We adopt the DAG formulation for three principled reasons.

\paragraph{Pragmatic Blueprint, Not Causal Claim.}
We emphasize that our DAG is a \emph{generative blueprint}: its edges encode statistical and probabilistic dependencies inferred from limited observations, not claims of ground-truth causality. This pragmatic interpretation aligns with the foundational Bayesian Network (BN) literature~\cite{pearl2009causality, koller2009probabilistic}, where DAGs serve as compact encodings of conditional independence relationships without necessarily asserting causal direction. The distinction is important: StructSynth uses the discovered DAG to \emph{organize} the generation process, not to make ontological claims about the data-generating mechanism.

\paragraph{Functional Necessity of Directionality.}
The structure-guided synthesis stage (Section~\ref{sec:synth}) generates feature values autoregressively: each feature is produced conditioned on values already generated for its parent nodes. This autoregressive process inherently requires a \emph{total ordering} over features to determine which values are available as conditioning context at each step. An undirected graph $\{A\text{---}B, B\text{---}C\}$ is ambiguous---it cannot specify whether $A$ should be generated before $C$ or vice versa. A DAG $\{A \to B \to C\}$ resolves this ambiguity by defining a topological order that directly translates into a layer-by-layer generation schedule. Even when the underlying statistical relationship between two features is symmetric (e.g., $\mathrm{Age} \leftrightarrow \mathrm{Income}$), the generation process \emph{must} select a direction---the DAG provides exactly this. The specific direction chosen does not affect \emph{whether} the dependency is captured, only the \emph{order} in which values are conditioned.

% \paragraph{Autoregressive Generation Demands Direction.}
% Consider the alternative: an undirected graph $\{A\text{---}B, B\text{---}C\}$. Without directionality, we cannot determine whether $A$ should be generated before $C$ or vice versa. Converting this to a DAG---either $A \to B \to C$ or $C \to B \to A$---resolves the ambiguity by defining a topological ordering that directly determines the generation schedule. Crucially, both orientations capture the same dependency between $A$ and $C$ (mediated through $B$); the choice of direction affects the conditioning order, not whether the dependency is respected.

\paragraph{Acyclicity as a Mathematical Prerequisite.}
Beyond directionality, acyclicity is essential for the generation process to be well-defined. If cycles were permitted---e.g., $A \to B \to C \to A$---the generation would encounter circular dependencies: the value of $A$ would depend on $C$, which depends on $B$, which in turn depends on $A$, creating an unresolvable loop. The DAG constraint guarantees the existence of a valid topological ordering over nodes, ensuring that every feature can be generated exactly once, conditioned only on features whose values have already been determined.

\paragraph{Summary.}
In summary, the DAG in \textsf{StructSynth} serves a \emph{functional} rather than \emph{ontological} role: it imposes a well-defined, acyclic generation order that enables structure-conditioned autoregressive synthesis. This formulation is consistent with decades of probabilistic graphical model research and is a natural fit for our discover-then-synthesize paradigm.


\subsection{Dependency Structure Discovery}

The first step in \textsf{StructSynth} is to uncover the dependency structure present in the limited training data, represented as a DAG $G = (V, E)$. The node set $V \subseteq \mathbf{A}$ consists of selected attributes from the dataset schema $\mathbf{A}$, while the edges $E$ capture probabilistic dependencies among these attributes. The framework employs an iterative, breadth-first search (BFS) guided by an LLM to progressively construct the dependency graph.

\paragraph{Design Rationale for Low-Data Regimes.}
A distinguishing aspect of this design is its explicit adaptation to data-scarce settings, achieved through three complementary mechanisms. \textbf{(a) LLM semantic prior as structural regularizer:} Rather than relying solely on statistical estimates from $n \le 100$ samples---which are often noisy and prone to spurious correlations---the framework leverages the LLM's world knowledge as an informative prior. This prior acts as a regularizer during structure discovery, filtering out statistically plausible but semantically implausible dependencies (e.g., a spurious correlation between \textit{Zip Code} and \textit{Income} that would mislead a purely data-driven algorithm). \textbf{(b) Statistical cues as weak supervision:} The computed association scores $\mathcal{S}$ serve as \emph{weak supervision} that grounds the LLM's semantic reasoning in empirical evidence, rather than serving as the sole decision criterion. This hybrid approach is crucial because neither source of evidence is reliable in isolation at small $n$: statistics are noisy, while LLM priors may not reflect dataset-specific idiosyncrasies. \textbf{(c) Decoupled architecture reduces per-stage sample complexity:} By separating structure discovery from data generation, each stage operates with a simpler objective that requires fewer samples to achieve reliably. Structure discovery only needs to identify which dependencies exist (a discrete, lower-dimensional problem), while generation only needs to model conditional distributions along the discovered graph---a factorization that is substantially easier than jointly learning both structure and generation parameters.

The core steps of the discovery process are described below.

\subsubsection{Graph Initialization}
We initiate the graph construction by identifying source nodes attributes that are not influenced by any others in the dataset. A small subset of the training data, $\mathcal{D}_{\mathtt{few\_shot}} \subseteq \mathcal{D}_{\mathtt{train}}$, serves as in-context examples for the LLM, which is guided by a specifically designed prompt, $\pi_{\mathtt{source}}$. Formally, the initial set of source nodes, $V_{\mathtt{source}}$, is determined by the following operation:
\begin{equation}
    V_{\mathtt{source}} = f_{\mathtt{LLM}}(\pi_{\mathtt{source}} ( \mathcal{D}_{\mathtt{few\_shot}})) .
\end{equation}
Consequently, the initial graph is set to $G_0 = (V_{\mathtt{source}}, \emptyset)$. To manage the discovery process, a queue $Q$ for the Breadth-First Search (BFS) is initialized with these source nodes, i.e., $Q \leftarrow V_{\mathtt{source}}$, and a set for tracking explored attributes, $V_{\mathtt{visited}}$, is initialized as empty.

\subsubsection{Expansion and Reasoned Link Generation}

The framework expands the graph $G$ through an iterative, breadth-first process. In each iteration $t$, a node $A_i$ is dequeued from the processing queue $Q$ and added to the set of visited nodes, $V_{\mathtt{visited}}$. To inform the model's link generation with empirical data, we compute a vector of association scores, $\mathcal{S}(A_i) \in \mathbb{R}^{K-1}$, between $A_i$ and all other attributes in $\mathbf{A}$. These scores are calculated using dependency measures specific to the data types of each attribute pair: \textit{Pearson's R} for continuous-continuous pairs, \textit{Correlation Ratio} for categorical-continuous pairs, and \textit{Cramér's V} for categorical-categorical pairs~\cite{agresti2011categorical} (detailed in Section~\ref{sec:association_scores}).

Subsequently, the LLM is queried with a comprehensive prompt, $\pi_{\mathtt{generate}}$, which integrates the current graph state $G_t$, the node $A_i$, and the statistical scores $\mathcal{S}(A_i)$. The model's task is to propose a set of direct successors for $A_i$ and provide a textual rationale for each proposed relationship. The structured output of this query is a set of successor-rationale pairs, denoted as $P_i$:
\begin{equation}
    P_i = f_{\mathtt{LLM}}(\pi_{\mathtt{generate}}(A_i, G_t, \mathcal{S}(A_i), \mathcal{D}_{\mathtt{few\_shot}})).
\end{equation}
From this set, we define the proposed new edges, $E_{\mathtt{prop}, t}$, and the set of newly discovered nodes, $V_{\mathtt{new}, t}$, which are not already in the graph's vertex set $V_t$:
\begin{equation}
    \begin{aligned}
        E_{\mathtt{prop}, t} &= \{ (A_i \to A_j, r_{ij}) \mid (A_j, r_{ij}) \in P_i \}, \\
        V_{\mathtt{new}, t} &= \{ A_j \mid (A_j, \_) \in P_i \land A_j \notin V_t \}.
    \end{aligned}
\end{equation}
These new nodes are then added to the queue $Q$ for subsequent exploration, ensuring the continued expansion of the graph.%  boundary.

\subsubsection{LLM-based Cycle Resolution}

To ensure that the graph remains a DAG throughout, the framework resolves any cycles that may result from integrating proposed edges $E_{\mathtt{prop}, t}$ into the current graph $G_t = (V_t, E_t)$. Specifically, we construct the candidate edge set $E_{\mathtt{candidate}} = E_t \cup E_{\mathtt{prop}, t}$ and apply a cycle detection algorithm to identify all elementary cycles, $\Psi$. 

If cycles are detected ($\Psi \neq \emptyset$), the framework initiates a reasoned resolution process. For each cycle $\psi \in \Psi$, where $\psi$ represents the specific subset of edges from $E_{\mathtt{candidate}}$ that forms the circular path (e.g., $\{ (A \to B, r_{ab}), (B \to C, r_{bc}), (C \to A, r_{ca}) \}$), the LLM is queried with the prompt $\pi_{\mathtt{resolve}}$. The prompt's input, $\psi$, provides the model with the complete logical contradiction, including the rationales for both newly proposed and previously established edges. The model's task is to analyze this conflict and identify the single edge whose removal is most justified. This resolution logic is applied to every detected cycle. The complete set of pruned edges, $E_{\mathtt{pruned}}$, is thus defined by collecting the output for each cycle in $\Psi$. The state for the next iteration, $G_{t+1}$, is then formally defined by removing these pruned edges from the candidate set:
\begin{equation}
    \begin{aligned}
        E_{\mathtt{pruned}} &= \{ f_{\mathtt{LLM}}(\pi_{\mathtt{resolve}}(\psi)) \mid \psi \in \Psi \} \\
        G_{t+1} &= (V_t \cup V_{\mathtt{new}, t}, (E_t \cup E_{\mathtt{prop}, t}) \setminus E_{\mathtt{pruned}}).
    \end{aligned}
\end{equation}
This general mechanism ensures the graph remains acyclic and empowers the model to revise its structure based on new, more confident evidence.


\subsubsection{Continuation of BFS}

The process continues iteratively. Following the cycle resolution step, the newly discovered nodes, $V_{\mathtt{new}, t}$, which are now part of the updated graph $G_{t+1}$, are added to the processing queue $Q$. The framework then returns to the \textit{Expansion and Reasoned Link Generation} step to process the next node from the queue. This iterative expansion-resolution cycle continues until $Q$ is empty, at which point the complete, acyclic dependency graph $G$ is returned as the final output.

\begin{table*}[ht]

    \small
    \centering
    \setlength{\tabcolsep}{1mm}
    \caption{Comparison of Models on Downstream Model Performance (\%). (C) and (R) denote classification and regression tasks, respectively. $^\dag$ indicates datasets created post LLM knowledge cutoff. Best results are in \textbf{bold}; second-best are \underline{underlined}.}
    \begin{tabular}{ll|llllll|c|c}
    \toprule
    
    \textbf{Type} & \textbf{Method} & 
    \multicolumn{1}{c}{\textbf{Adult (C)}} & 
    \multicolumn{1}{c}{\textbf{Anxiety (C)$^\dag$}} & 
    \multicolumn{1}{c}{\textbf{Compas (C)}} & 
    \multicolumn{1}{c}{\textbf{Salary (R)$^\dag$}} & 
    \multicolumn{1}{c}{\textbf{Obesity (R)}} & 
    \multicolumn{1}{c}{\textbf{Churn (C)}} & 
    \textbf{Average} & \textbf{Avg. Rank} \\
    \midrule
    \multirow{1}{*}{Real} & $\mathcal{D}_\mathtt{train}$ & $82.51_{\pm 1.93}$ & $85.49_{\pm 1.60}$ & $64.84_{\pm 2.51}$ & $49.71_{\pm 7.02}$ & $59.83_{\pm 6.40}$ & $86.86_{\pm 3.78}$ & $71.54$ & - \\
    
    \midrule
    \multirow{5}{*}{DGMs} & TVAE & $81.21_{\pm 1.54}$ & $79.63_{\pm 2.65}$ & $64.39_{\pm 2.06}$ & $46.18_{\pm 5.91}$ & $60.92_{\pm 2.61}$ & $87.42_{\pm 3.80}$ & $69.96$ & $6.50$ \\
    & DDPM & $81.72_{\pm 1.79}$ & $82.94_{\pm 3.24}$ & $66.01_{\pm 2.50}$ & $22.30_{\pm 15.70}$ & $43.20_{\pm 7.98}$ & $85.91_{\pm 2.05}$ & $63.68$ & $8.17$ \\
    & CTGAN & $81.14_{\pm 1.81}$ & $76.06_{\pm 1.47}$ & $65.17_{\pm 2.61}$ & $46.17_{\pm 5.25}$ & $54.29_{\pm 5.90}$ & $86.50_{\pm 4.17}$ & $68.22$ & $7.83$ \\
    & NFlow & $75.42_{\pm 4.58}$ & $74.83_{\pm 2.94}$ & $63.68_{\pm 4.22}$ & $25.55_{\pm 6.16}$ & $46.80_{\pm 4.17}$ & $84.99_{\pm 2.73}$ & $61.88$ & $11.00$ \\
    & TabSyn & $83.88_{\pm 1.56}$ & \underline{$85.68_{\pm 2.40}$} & $67.77_{\pm 6.98}$ & $47.17_{\pm 9.58}$ & $43.50_{\pm 9.91}$ & $89.86_{\pm 1.53}$ & $69.64$ & $3.83$ \\
    \midrule
    \multirow{3}{*}{\makecell[l]{Structure\\Aware}} & GOGGLE & $78.92_{\pm 3.04}$ & $80.92_{\pm 1.61}$ & $63.73_{\pm 3.85}$ & $43.18_{\pm 10.65}$ & $55.84_{\pm 7.77}$ & $85.27_{\pm 3.98}$ & $67.98$ & $8.50$ \\
    & BN & $82.68_{\pm 1.35}$ & $82.77_{\pm 2.39}$ & $60.13_{\pm 3.46}$ & $46.33_{\pm 7.29}$ & $59.76_{\pm 3.46}$ & $88.85_{\pm 3.02}$ & $70.09$ & $6.00$ \\
    & DECAF & $77.64_{\pm 3.91}$ & $79.90_{\pm 2.30}$ & $63.90_{\pm 3.69}$ & $43.58_{\pm 6.25}$ & $52.88_{\pm 6.44}$ & $83.94_{\pm 2.97}$ & $66.97$ & $9.17$ \\
    \midrule
    \multirow{3}{*}{LLMs} & GReaT & $82.55_{\pm 2.20}$ & $83.81_{\pm 2.55}$ & $64.02_{\pm 2.72}$ & $51.43_{\pm 3.31}$ & $47.03_{\pm 9.16}$ & $86.67_{\pm 3.20}$ & $69.25$ & $6.00$ \\
    & CLLM & \underline{$83.95_{\pm 2.09}$} & $85.17_{\pm 0.94}$ & \underline{$66.24_{\pm 1.51}$} & \underline{$54.53_{\pm 0.20}$} & \underline{$60.43_{\pm 7.19}$} & \underline{$89.81_{\pm 2.78}$} & \underline{$73.36$} & \underline{$2.67$} \\
    & GraDe & $83.72_{\pm 1.44}$ & $79.74_{\pm 2.01}$ & $66.15_{\pm 4.90}$ & $35.03_{\pm 8.08}$ & $22.92_{\pm 7.01}$ & $88.81_{\pm 2.51}$ & $62.73$ & $7.33$ \\
    \midrule
    \multirow{1}{*}{Ours} & StructSynth & $\boldsymbol{85.55_{\pm 1.08}}$ & $\boldsymbol{86.45_{\pm 0.57}}$ & $\boldsymbol{69.40_{\pm 0.25}}$ & $\boldsymbol{55.98_{\pm 2.96}}$ & $\boldsymbol{62.27_{\pm 5.40}}$ & $\boldsymbol{90.39_{\pm 1.96}}$ & $\boldsymbol{75.01}$ & $\boldsymbol{1.00}$ \\
    
    \bottomrule
    \end{tabular}
    \label{tab:results_comparison_performance}
\end{table*}

\subsection{Structure-Guided Synthesis}
\label{sec:synth}

Upon learning the dependency structure $G=(V, E)$, we utilize it to guide synthetic tabular data generation. Constructing each synthetic data point $\tilde{\mathbf{x}}_j$ is performed in two sequential phases aligned with the underlying data structure.

\subsubsection{Generating Graph-Based Values}

This phase generates values for attributes $V$ defined by the learned dependency graph. Initially, the nodes $V$ are partitioned into topological layers $\mathcal{L} = (L_1, \dots, L_m)$. For each layer $L_i$, we extract the corresponding dependency subgraph $G_i = G[L_i \cup \Gamma^{-}(L_i)]$, where $\Gamma^{-}(L_i)$ represents the parent nodes of attributes in $L_i$. Additionally, we select the relevant slice of the few-shot dataset as $\mathcal{D}_{\mathtt{few\_shot}}^i = \Pi_{L_i \cup \Gamma^{-}(L_i)}(\mathcal{D}_{\mathtt{few\_shot}})
$, where $\Pi$ denotes the projection operator onto specified columns.

Values for attributes in layer $L_i$, denoted as $\tilde{\mathbf{x}}_{j, L_i}$, are generated conditionally based on previously generated values for preceding layers $\tilde{\mathbf{x}}_{j, <i}$:
\begin{equation}
\tilde{\mathbf{x}}_{j, L_i} = f_{\mathtt{LLM}}\left(\pi_{\mathtt{data\_gen}}\left( \tilde{\mathbf{x}}_{j, <i},\ G_i,\ L_i,\ \mathcal{D}_{\mathtt{few\_shot}}^i\right)\right).
\end{equation}
Iterating through all layers yields a partial synthetic data point $\tilde{\mathbf{x}}_{j, V}$ comprising values for all graph-dependent attributes.

\subsubsection{Generating Independent Values}

Subsequently, values for isolated attributes $A_{\mathtt{iso}} = \mathbf{A} \setminus V$ are generated collectively in a single step, conditioned upon the complete set of previously generated graph-based values $\tilde{\mathbf{x}}_{j, V}$:
\begin{equation}
\tilde{\mathbf{x}}_{j, A_{\mathtt{iso}}} = f_{\mathtt{LLM}}\left(\pi_{\mathtt{iso}}\left(\tilde{\mathbf{x}}_{j, V},\ A_{\mathtt{iso}},\ \mathcal{D}_{\mathtt{few\_shot}}^{\mathtt{iso}}\right)\right),
\end{equation}
where $\mathcal{D}_{\mathtt{few\_shot}}^{\mathtt{iso}} = \Pi_{A_{\mathtt{iso}}}(\mathcal{D}_{\mathtt{few\_shot}}).$

\subsubsection{Final Assembly}

A complete synthetic data point $\tilde{\mathbf{x}}_j$ is formed by concatenating the graph-dependent and independent attribute values as $\tilde{\mathbf{x}}_j = \mathtt{concat}(\tilde{\mathbf{x}}_{j, V}, \tilde{\mathbf{x}}_{j, A_{\mathtt{iso}}})$. Repeating this procedure $s$ times yields the synthetic dataset $\mathbf{X}_{\mathtt{synth}} = \{\tilde{\mathbf{x}}_1, \ldots, \tilde{\mathbf{x}}_s\}$, resulting in the final synthetic dataset $\mathcal{D}_{\mathtt{synth}} = (\mathbf{X}_{\mathtt{synth}}, \mathbf{A})$.

\subsection{Statistical Association Measures}
\label{sec:association_scores}

To quantify the empirical dependencies between attribute pairs, the framework employs three association measures, chosen according to the data types of the variables involved.

\subsubsection{Continuous-Continuous: Pearson's R}
For two continuous variables $x = \{x_i\}_{i=1}^n$ and $y = \{y_i\}_{i=1}^n$, Pearson's correlation coefficient $r$ is defined as the standardized covariance:
\begin{equation}
    r = \frac{\sum_{i=1}^{n}(x_i-\bar{x})(y_i-\bar{y})}{\sqrt{\sum_{i=1}^{n}(x_i-\bar{x})^2}\sqrt{\sum_{i=1}^{n}(y_i-\bar{y})^2}}
\end{equation}
where $\bar{x}$ and $\bar{y}$ denote the sample means. The value of $r$ ranges from $[-1, 1]$, where the sign indicates the direction of the linear relationship and the absolute value indicates its strength.

\subsubsection{Categorical-Continuous: Correlation Ratio}
The Correlation Ratio, $\eta$, measures the relationship between a categorical variable $x$ and a continuous variable $y$. It is defined as the square root of the weighted variance of the category means divided by the total variance:
\begin{equation}
    \eta = \sqrt{\frac{\sum_{x} n_x(\bar{y}_x - \bar{y})^2}{\sum_{i=1}^{n}(y_i - \bar{y})^2}}
\end{equation}
where $n_x$ is the number of samples in category $x$, $\bar{y}_x$ is the mean of $y$ for samples in category $x$, and $\bar{y}$ is the overall mean. The value of $\eta$ ranges from $[0, 1]$, with values closer to 1 indicating that the categorical variable provides more information about the continuous variable.

\subsubsection{Categorical-Categorical: Cramér's V}
For two categorical variables with a contingency table of dimensions $r \times k$, Cramér's V is based on the $\chi^2$ statistic. To ensure robustness, we employ the bias correction proposed by Bergsma and Wicher (2013).

First, the bias-corrected $\phi^2$ is calculated as:
\begin{equation}
    \phi^2_{\text{corr}} = \max\left(0, \frac{\chi^2}{n} - \frac{(k-1)(r-1)}{n-1}\right)
\end{equation}
where $n$ is the total number of samples. The corrected dimensions are:
\begin{equation}
    r_{\text{corr}} = r - \frac{(r-1)^2}{n-1}, \quad k_{\text{corr}} = k - \frac{(k-1)^2}{n-1}
\end{equation}
The final Cramér's V score is then given by:
\begin{equation}
    V_{\text{corr}} = \sqrt{\frac{\phi^2_{\text{corr}}}{\min(k_{\text{corr}}-1, r_{\text{corr}}-1)}}
\end{equation}
The value ranges from $[0, 1]$, representing the strength of association between the two categorical variables.

\subsection{Complete Algorithm}
\label{sec:algorithm}

Algorithm~\ref{alg:structsynth} provides the complete pseudocode for the \textsf{StructSynth} framework, integrating both the dependency structure discovery and the structure-guided synthesis stages described above.

\begin{algorithm*}
\caption{StructSynth: Structure-Aware Tabular Data Synthesis}
\label{alg:structsynth}
\begin{algorithmic}[1]
\State \textbf{Input:} Training data $\mathcal{D}_{\mathtt{train}} = (X_{\mathtt{train}}, A)$, number of synthetic samples $s$
\State \textbf{Output:} Synthetic dataset $\mathcal{D}_{\mathtt{synth}} = (X_{\mathtt{synth}}, A)$

\vspace{0.2cm}
\Statex \textbf{Stage 1: Dependency Structure Discovery}
\Statex \textit{1. Graph Initialization}
\State Select a few-shot subset $\mathcal{D}_{\mathtt{few\_shot}} \subseteq \mathcal{D}_{\mathtt{train}}$
\State $V_{\mathtt{source}} \gets f_{\mathtt{LLM}}(\pi_{\mathtt{source}}(\mathcal{D}_{\mathtt{few\_shot}}))$
\State $G \gets (V_{\mathtt{source}}, \emptyset)$
\State $Q \gets V_{\mathtt{source}}$ \Comment{Initialize queue for BFS}
\State $V_{\mathtt{visited}} \gets \emptyset$

\Statex \textit{2. Iterative Expansion, Resolution, and Continuation (BFS)}
\While{$Q \neq \emptyset$}
    \Statex \hspace{\algorithmicindent} \textit{a. Expansion and Reasoned Link Generation}
    \State $A_i \gets \mathtt{Dequeue}(Q)$
    \State $V_{\mathtt{visited}} \gets V_{\mathtt{visited}} \cup \{A_i\}$
    \State Compute association scores $\mathcal{S}(A_i)$ between $A_i$ and other attributes
    \State $P_i \gets f_{\mathtt{LLM}}(\pi_{\mathtt{generate}}(A_i, G, \mathcal{S}(A_i), \mathcal{D}_{\mathtt{few\_shot}}))$
    \State $E_{\mathtt{prop},t} \gets \{(A_i \to A_j, r_{ij}) \mid (A_j, r_{ij}) \in P_i\}$
    \State $V_{\mathtt{new},t} \gets \{A_j \mid (A_j, \_) \in P_i \land A_j \notin V\}$
    
    \Statex \hspace{\algorithmicindent} \textit{b. LLM-based Cycle Resolution}
    \State $E_{\mathtt{candidate}} \gets E \cup E_{\mathtt{prop},t}$
    \State $\Psi \gets \mathtt{DetectCycles}(V \cup V_{\mathtt{new},t}, E_{\mathtt{candidate}})$
    \State $E_{\mathtt{pruned}} \gets \emptyset$
    \For{each cycle $\psi \in \Psi$}
        \State $e_{\mathtt{to\_remove}} \gets f_{\mathtt{LLM}}(\pi_{\mathtt{resolve}}(\psi))$
        \State $E_{\mathtt{pruned}} \gets E_{\mathtt{pruned}} \cup \{e_{\mathtt{to\_remove}}\}$
    \EndFor
    
    \Statex \hspace{\algorithmicindent} \textit{c. Continuation of BFS}
    \State $G \gets (V \cup V_{\mathtt{new},t}, E_{\mathtt{candidate}} \setminus E_{\mathtt{pruned}})$ \Comment{Update graph}
    \State Enqueue unvisited nodes from $V_{\mathtt{new},t}$ into $Q$
\EndWhile
\State \textbf{Return} $G = (V, E)$

\vspace{0.3cm}

\Statex \textbf{Stage 2: Structure-Guided Synthesis}
\State Partition nodes $V$ into topological layers $\mathcal{L} = (L_1, \dots, L_m)$
\State $X_{\mathtt{synth}} \gets \emptyset$
\For{$j \gets 1$ to $s$}
    \Statex \hspace{\algorithmicindent} \textit{a. Generating Graph-Based Values}
    \State $\tilde{x}_{j,V} \gets \{\}$
    \For{$i \gets 1$ to $m$}
        \State $L_i \gets \mathcal{L}[i]$
        \State $\tilde{x}_{j,L_i} \gets f_{\mathtt{LLM}}\Big(\pi_{\mathtt{data\_gen}}(\tilde{x}_{j,<i}, G[L_i \cup \Gamma^-(L_i)], L_i, \mathcal{D}_{\mathtt{few\_shot}}^i)\Big)$
        \State $\tilde{x}_{j,V} \gets \tilde{x}_{j,V} \cup \tilde{x}_{j,L_i}$
    \EndFor
    
    \Statex \hspace{\algorithmicindent} \textit{b. Generating Independent Values}
    \State $A_{\mathtt{iso}} \gets A \setminus V$
    \State $\tilde{x}_{j, A_{\mathtt{iso}}} \gets f_{\mathtt{LLM}}(\pi_{\mathtt{iso}}(\tilde{x}_{j,V}, A_{\mathtt{iso}}, \mathcal{D}_{\mathtt{few\_shot}}^{\mathtt{iso}}))$
    
    \Statex \hspace{\algorithmicindent} \textit{c. Final Assembly}
    \State $\tilde{x}_j \gets \mathtt{concat}(\tilde{x}_{j,V}, \tilde{x}_{j,A_{\mathtt{iso}}})$
    \State $X_{\mathtt{synth}} \gets X_{\mathtt{synth}} \cup \{\tilde{x}_j\}$
\EndFor
\State \textbf{Return} $\mathcal{D}_{\mathtt{synth}} = (X_{\mathtt{synth}}, A)$
\end{algorithmic}
\end{algorithm*}
